\documentclass[border=12pt]{standalone}
\usepackage[utf8]{inputenc}
\usepackage[T1]{fontenc}
\usepackage[american, siunitx]{circuitikz}

\sisetup{output-decimal-marker={,}}

\ctikzset{
    resistors/scale=0.8
}

\begin{document}
\begin{circuitikz}[font=\sffamily\small]

    % =========================================================================
    % 1. BLOCO ESQUERDO: ARDUINO UNO R3 - MESTRE (U1)
    % =========================================================================
    \draw[thick, fill=gray!5] (-6.0, -3.0) rectangle (-1.2, 3.0);
    \node[align=center, font=\sffamily\bfseries] at (-4.2, 0.4) {U1\\Arduino\\Uno R3};
    \node[font=\sffamily\scriptsize, text=gray!70!black] at (-4.2, -0.7) {(Mestre)};

    % Alimentacao de U1
    \draw (-3.6, 3.0) -- (-3.6, 3.8) node[vcc]{+5V};
    \node[below, font=\ttfamily\footnotesize] at (-3.6, 3.0) {5V};

    \draw (-3.6, -3.0) -- (-3.6, -3.8) node[ground]{};
    \node[above, font=\ttfamily\footnotesize] at (-3.6, -3.0) {GND};

    % Pinos I2C de U1
    \node[left, font=\ttfamily\scriptsize] at (-1.2, 0.8) {A4 (SDA)};
    \node[left, font=\ttfamily\scriptsize] at (-1.2, -0.8) {A5 (SCL)};

    % =========================================================================
    % 2. BARRAMENTO I2C COMPARTILHADO E RESISTORES DE PULL-UP (R1, R2 = 10k)
    % =========================================================================
    % Linhas principais SDA e SCL (sem setas nos nos de juncao em T)
    \draw[<-, >=stealth, thick] (-1.2, 0.8) -- (4.2, 0.8) node[circ]{};
    \node[above, font=\ttfamily\scriptsize] at (-0.3, 0.8) {SDA};

    \draw[thick] (-1.2, -0.8) -- (5.0, -0.8) node[circ]{};
    \node[below, font=\ttfamily\scriptsize] at (-0.3, -0.8) {SCL};

    % Barramento +5V dos resistores de pull-up (ponto apenas na juncao em T)
    \draw (0.4, 3.6) -- (2.6, 3.6);
    \node[vcc] at (1.5, 3.6) {+5V};
    \node[circ] at (1.5, 3.6) {};

    % Pull-up R1 (10 kOhm) na linha SDA
    \draw (0.4, 0.8) node[circ]{} -- (0.4, 1.4)
        to[R, l^=R1, a_=10<\kilo\ohm>] (0.4, 3.6);

    % Pull-up R2 (10 kOhm) na linha SCL
    \draw (2.6, -0.8) node[circ]{} -- (2.6, 1.4)
        to[R, l^=R2, a_=10<\kilo\ohm>] (2.6, 3.6);

    % Derivacoes verticais para os escravos (setas apenas nos pinos de entrada)
    % Tronco SDA -> LCD1 e U2
    \draw[->, >=stealth, thick] (4.2, 0.8) -- (4.2, 4.2) -- (6.2, 4.2);
    \draw[->, >=stealth, thick] (4.2, 0.8) -- (4.2, -2.8) -- (6.2, -2.8);

    % Tronco SCL -> LCD1 e U2
    \draw[->, >=stealth, thick] (5.0, -0.8) -- (5.0, 2.8) -- (6.2, 2.8);
    \draw[->, >=stealth, thick] (5.0, -0.8) -- (5.0, -4.2) -- (6.2, -4.2);

    % =========================================================================
    % 3. ESCRAVO SUPERIOR: DISPLAY LCD 20x4 I2C (LCD1 - 0x27)
    % =========================================================================
    \draw[thick, fill=gray!5] (6.2, 1.8) rectangle (11.4, 5.6);
    \node[align=center, font=\sffamily\bfseries] at (9.1, 4.1) {LCD1\\Display LCD\\20$\times$4};
    \node[font=\sffamily\scriptsize, text=gray!70!black] at (9.1, 2.8) {(I2C -- 0x27)};

    % Alimentacao de LCD1
    \draw (8.8, 5.6) -- (8.8, 6.4) node[vcc]{+5V};
    \node[below, font=\ttfamily\scriptsize] at (8.8, 5.6) {VCC};

    \draw (8.8, 1.8) -- (8.8, 1.0) node[ground]{};
    \node[above, font=\ttfamily\scriptsize] at (8.8, 1.8) {GND};

    % Pinos I2C de LCD1
    \node[right, font=\ttfamily\scriptsize] at (6.2, 4.2) {SDA};
    \node[right, font=\ttfamily\scriptsize] at (6.2, 2.8) {SCL};

    % =========================================================================
    % 4. ESCRAVO INFERIOR: SENSOR MPU-6050 (U2 - 0x68)
    % =========================================================================
    \draw[thick, fill=gray!5] (6.2, -5.6) rectangle (11.4, -1.8);
    \node[align=center, font=\sffamily\bfseries] at (9.1, -3.4) {U2\\MPU-6050};
    \node[font=\sffamily\scriptsize, text=gray!70!black] at (9.1, -4.3) {(I2C -- 0x68)};

    % Alimentacao de U2
    \draw (8.8, -1.8) -- (8.8, -1.0) node[vcc]{+5V};
    \node[below, font=\ttfamily\scriptsize] at (8.8, -1.8) {VCC};

    \draw (8.8, -5.6) -- (8.8, -6.4) node[ground]{};
    \node[above, font=\ttfamily\scriptsize] at (8.8, -5.6) {GND};

    % Pinos I2C de U2
    \node[right, font=\ttfamily\scriptsize] at (6.2, -2.8) {SDA};
    \node[right, font=\ttfamily\scriptsize] at (6.2, -4.2) {SCL};

\end{circuitikz}
\end{document}